%%%%%%%%%%%%%%%%%%%%%%%%%%%%%%%%%%%%%%%%%%%%%%%%%%%%%%%%%%%%%%%%%%%%%%%%%%%
%%  chapter3.tex  –  Chapter 3: Adaptation and Alignment to Network
%%                   Intrusion Detection and Prevention Systems
%%  v9 – Domain-specific chapter: terminologies, data→graph transform,
%%       taxonomy (34 classes), features (83), hyperparameters,
%%       method instantiations. No experimental results (see Chapter 4).
%%       Only \section headings go to TOC.
%%%%%%%%%%%%%%%%%%%%%%%%%%%%%%%%%%%%%%%%%%%%%%%%%%%%%%%%%%%%%%%%%%%%%%%%%%%

\chapter{Adaptation and Alignment to Network Intrusion Detection and Prevention Systems}
\label{ch:adaptation}

This chapter maps the domain-independent mathematical methods developed in Chapter~\ref{ch:methods} onto the entities and constraints of cybersecurity and, specifically, hybrid Network Intrusion Detection and Prevention Systems (NIDPS). The abstract objects of the preceding chapter --- temporal graphs, state-space dynamics, variational posteriors, transport plans, and game-theoretic utilities --- are instantiated with the concrete entities of network security: flows, packets, sessions, hosts, services, and adversarial actors. The chapter defines the cybersecurity-specific terminologies, describes how raw network traffic is transformed into graph structures suitable for input into the developed GNN architectures, specifies the classification taxonomy and feature engineering pipeline, presents the hyperparameter configuration and optimisation rationale for each model, and demonstrates how each method from Chapter~\ref{ch:methods} is adapted with domain-specific modifications. No experimental results are presented here; these are reserved for Chapter~\ref{ch:experiments}.


%%=====================================================================
\section{Network Intrusion Detection: Terminologies and Context}
\label{sec:nidps_context}
%%=====================================================================

\subsection{Hybrid Intrusion Detection and Prevention Systems}

A hybrid intrusion detection and prevention system (HIDPS)~\cite{Buczak2016,Ahmad2021,Aldweesh2020,Ferrag2020,Paxson1999} combines two complementary detection paradigms. Signature-based detection matches incoming traffic against a database of known patterns (rules) and achieves high precision on previously catalogued threats but cannot detect novel attacks. Anomaly-based detection learns a statistical model of normal behaviour and flags deviations, enabling detection of zero-day exploits and evolving attack patterns at the cost of potentially higher false-positive rates. The dissertation's methods operate within the anomaly-based branch, augmented by large language model reasoning for zero-shot classification of novel categories.

A Network IDPS (NIDPS) monitors traffic at the network layer, analysing flows and packets between hosts rather than events within a single endpoint. The three operating conditions of the architectural framework (Figure~\ref{fig:dissertation_architecture}) are all present in this domain: adversarial manipulation by deliberate opponents who craft traffic to evade detection, non-stationarity through evolving traffic patterns and protocol changes, and distribution of the monitoring infrastructure across organisational boundaries.


\subsection{Key Terminologies}

The following terms are used throughout this chapter and the experimental validation in Chapter~\ref{ch:experiments}:

\noindent\textbf{Flow.} A sequence of packets sharing a common 5-tuple (source IP, destination IP, source port, destination port, protocol). A flow is the fundamental unit of analysis and corresponds to a directed edge in the network graph.

\noindent\textbf{Session.} A bidirectional communication comprising a forward flow and its reverse counterpart, capturing the full exchange between two endpoints.

\noindent\textbf{Host.} A network endpoint (IP address) that participates as either source or destination. Hosts correspond to nodes in the network graph.

\noindent\textbf{Service.} A functional component identified by port number and protocol (e.g., HTTP on port 80, SSH on port 22). In cloud-native deployments, services correspond to microservice containers.

\noindent\textbf{PCAP (Packet Capture).} The raw binary format capturing full packet headers and optionally payloads, from which flow-level features are extracted.

\noindent\textbf{IDS Rule.} A structured pattern specification (e.g., Snort or Suricata rule syntax) encoding a detection signature with action, protocol, addresses, ports, and content matching directives.

\noindent\textbf{MITRE ATT\&CK.} A comprehensive knowledge base of adversary tactics, techniques, and procedures (TTPs) used to classify and contextualise detected threats.

\noindent\textbf{SOC (Security Operations Centre).} The organisational unit responsible for monitoring, analysing, and responding to security events, staffed by analysts at Tiers~1--3.


%%=====================================================================
\section{Data-to-Graph Transformation Pipeline}
\label{sec:data_to_graph}
%%=====================================================================

The central premise of this dissertation is that network traffic is naturally graph-structured: hosts are nodes, flows are edges, and temporal ordering defines the graph evolution. This section formalises the transformation from raw network data to the graph inputs required by the developed GNN architectures.

\subsection{Graph Construction from Network Traffic}

Given a stream of network flows $\{(s_i, d_i, t_i, \mathbf{f}_i)\}_{i=1}^{N}$, where $s_i$ is the source host, $d_i$ is the destination host, $t_i$ is the timestamp, and $\mathbf{f}_i\in\mathbb{R}^{83}$ is the feature vector, the temporal graph $\calG(t)=(\calV,\calE(t),\mathbf{X}(t))$ is constructed as follows:

\noindent\textbf{Nodes.} Each unique IP address becomes a node $v\in\calV$. Node features $\bx_v(t)$ aggregate the host-level statistics computed over a sliding window $[t-\Delta_w, t]$: cumulative byte counts (sent/received), active connection counts, unique destination ports contacted, protocol distribution (TCP/UDP/ICMP ratios), and mean inter-connection interval. This yields a 12-dimensional node feature vector.

\noindent\textbf{Edges.} Each flow $(s_i, d_i, t_i, \mathbf{f}_i)$ creates a directed edge from node $s_i$ to node $d_i$ at time $t_i$ with edge features $\mathbf{e}_{s_id_i}(t_i) = \mathbf{f}_i$. The full 83-dimensional feature vector (Section~\ref{sec:features}) is attached to the edge.

\noindent\textbf{Temporal evolution.} The graph evolves as new flows arrive and old connections expire. A flow is considered active if its most recent packet arrived within a configurable timeout $\Delta_{\mathrm{timeout}}=120$\,s. Expired edges are removed, yielding a dynamic, time-varying topology.

\subsection{Multi-Level Graph Construction}

For the TripleE-TGNN (Method M3) and the SOC Intelligence integration, the network traffic is additionally mapped to three hierarchical levels:

\noindent\textbf{Service graph $\calG_s(t)$.} Nodes represent services (port/protocol pairs); edges represent API calls between services, aggregated over configurable windows. Used for detecting inter-service attack patterns in microservices environments.

\noindent\textbf{Trace graph $\calG_r(t)$.} Nodes represent distributed request traces (identified by trace IDs in HTTP headers or OpenTelemetry instrumentation); edges connect traces that share at least one service. Used for detecting multi-stage attack chains.

\noindent\textbf{Host graph $\calG_n(t)$.} Nodes represent infrastructure hosts; edges represent observed network flows. This is the primary graph for all methods.

\noindent Bipartite coupling graphs $\calB_{sr}$, $\calB_{sn}$, $\calB_{rn}$ link entities across levels, enabling the cross-level attention mechanism of Equation~\eqref{eq:triple_cross_attn}.


%%=====================================================================
\section{Classification Taxonomy}
\label{sec:taxonomy}
%%=====================================================================

The platform classifies network traffic into 34 distinct classes: 1 benign class and 33 attack classes organised into 7 major categories. This taxonomy is unified across all six benchmark datasets, with dataset-specific attack types mapped to the canonical categories. Table~\ref{tab:taxonomy} presents the complete taxonomy.

\begin{table}[htbp]
\centering
\caption{Unified classification taxonomy: 34 classes across 7 attack categories.}
\label{tab:taxonomy}
\small
\begin{tabular}{p{2.5cm} p{10.5cm}}
\toprule
\textbf{Category} & \textbf{Attack Classes} \\
\midrule
\textbf{DDoS} (8) & DDoS-TCP, DDoS-UDP, DDoS-ICMP, DDoS-HTTP, DDoS-SYN Flood, DDoS-DNS Amplification, DDoS-SSDP, DDoS-NTP Amplification \\
\addlinespace
\textbf{DoS} (5) & DoS GoldenEye, DoS Hulk, DoS Slowhttptest, DoS Slowloris, DoS HeartBeat \\
\addlinespace
\textbf{Reconnaissance} (4) & PortScan, Nmap Scan, Service Probe, OS Fingerprint \\
\addlinespace
\textbf{Brute Force} (4) & FTP-Patator, SSH-Patator, Web Login Brute Force, RDP Brute Force \\
\addlinespace
\textbf{Web Attacks} (4) & XSS, SQL Injection, Command Injection, Directory Traversal \\
\addlinespace
\textbf{Malware/Botnet} (5) & Mirai, Botnet Communication, Trojan Callback, Ransomware Beacon, Cryptomining \\
\addlinespace
\textbf{Advanced} (3) & Infiltration, Heartbleed, Spoofing \\
\addlinespace
\textbf{Benign} (1) & Normal traffic \\
\bottomrule
\end{tabular}
\end{table}

The taxonomy is extensible: new attack classes can be registered through the platform administration interface without retraining the full model ensemble, leveraging the CyberSecLLM zero-shot classification capability and the TripleE-TGNN elastic weight consolidation for incremental learning.


%%=====================================================================
\section{Feature Engineering Pipeline}
\label{sec:features}
%%=====================================================================

All models operate on a shared vector of 83 engineered features extracted from raw PCAP files or pre-processed CSV datasets. Features span five categories, summarised in Table~\ref{tab:features}.

\begin{table}[htbp]
\centering
\caption{The 83-feature engineering pipeline organised by category.}
\label{tab:features}
\small
\begin{tabular}{p{3.5cm} p{1.5cm} p{8.5cm}}
\toprule
\textbf{Category} & \textbf{Count} & \textbf{Representative Features} \\
\midrule
Flow-level statistics & 22 & Duration, forward/backward packet counts, total bytes (fwd/bwd), packet length (mean, max, min, std), flow bytes/s, flow packets/s, flow IAT (mean, std, max, min) \\
\addlinespace
Inter-arrival time distributions & 18 & Forward IAT (mean, std, min, max), backward IAT (mean, std, min, max), active/idle time statistics, sub-flow forward/backward counts \\
\addlinespace
Flag-based indicators & 15 & SYN count, FIN count, RST count, PSH count, ACK count, URG count, CWE count, ECE count, forward/backward PSH flags, forward/backward URG flags, FIN flag count, SYN flag count, RST flag count \\
\addlinespace
Payload statistics & 16 & Segment sizes (min, max, mean), header lengths (mean, total), initial window bytes (forward/backward), payload byte entropy, packet length variance, average packet size, average segment size \\
\addlinespace
Derived ratios & 12 & Bytes-per-packet (fwd/bwd), packets-per-second, down/up ratio, average packet size ratio, flow IAT ratio, bulk rate (forward/backward), active mean/std, idle mean/std \\
\midrule
\textbf{Total} & \textbf{83} & \\
\bottomrule
\end{tabular}
\end{table}

Feature extraction is performed by a shared \texttt{FeatureExtractor} module using per-feature z-score standardisation fitted on the training distribution:
\begin{equation}
\tilde{f}_j = \frac{f_j - \mu_j^{\mathrm{train}}}{\sigma_j^{\mathrm{train}} + \varepsilon},
\label{eq:zscore}
\end{equation}
where $\mu_j^{\mathrm{train}}$ and $\sigma_j^{\mathrm{train}}$ are the mean and standard deviation of feature $j$ computed on the training set, and $\varepsilon=10^{-8}$ prevents division by zero. Datasets with fewer than 83 native features are zero-padded to the standard dimension; those with more are projected via PCA retaining 95\% explained variance.


%%=====================================================================
\section{Domain-Specific Instantiation of Methods}
\label{sec:method_instantiation}
%%=====================================================================

This section describes how each of the seven methods from Chapter~\ref{ch:methods} is instantiated with cybersecurity-specific parameters, demonstrating the mapping from abstract mathematical objects to concrete network-security entities.


\subsection{M1: CT-TGNN Applied to Network-Traffic Graph Dynamics}

CT-TGNN~\cite{Anaedevha2025cais} is instantiated as follows. In the network security instantiation, the node set $\calV$ comprises hosts and services, and the edge set $\calE(t)$ records flows. Node features $\bx_v(t)$ encode host-level statistics (12 dimensions). Edge features $\mathbf{e}_{uv}(t)$ capture the 83-dimensional flow feature vector from Table~\ref{tab:features}.

The three parallel ODE branches of Equation~\eqref{eq:ct_tgnn_multi} are assigned time constants $\tau_1=1$\,s (packet-level transients), $\tau_2=60$\,s (flow-level session dynamics), and $\tau_3=3600$\,s (campaign-level trends). This enables simultaneous resolution of sub-second port-scanning bursts and multi-hour advanced persistent threat campaigns. The architecture employs four message-passing layers, hidden dimension 256, state dimension 64, and DOPRI5 adaptive integration with tolerances $10^{-5}$ (relative) and $10^{-7}$ (absolute).

\begin{figure}[H]
\centering
\includegraphics[width=0.69\textwidth]{figures/fig_CT-TGNN_architecture.pdf}
\caption{CT-TGNN architecture instantiated for network traffic dynamics.}
\label{fig:ct_tgnn_arch}
\end{figure}


\subsection{M2: FedLLM-API for Collaborative Threat Intelligence}

FedLLM-API~\cite{Anaedevha2026neurips} is instantiated for collaborative threat intelligence. Each federated client corresponds to a security operations centre possessing local traffic datasets. Byzantine-robust aggregation (Algorithm~\ref{alg:fedllm}) protects against compromised participants submitting poisoned gradient updates. Differential privacy through local gradient clipping at threshold $C=1.0$ and Gaussian noise with scale $\sigma=0.8$ achieves an $(\epsilon{=}0.85,\delta{=}10^{-5})$ guarantee over 100 communication rounds. Sinkhorn OT alignment with regularisation $\lambda=0.1$ corrects for distributional heterogeneity across centres. The LLM ensemble weight is $\lambda_{\mathrm{LLM}}=0.3$.

\begin{figure}[H]
\centering
\includegraphics[width=0.96\textwidth]{figures/fig_fedllm_api_architecture.pdf}
\caption{FedLLM-API architecture for collaborative threat intelligence across organisational boundaries.}
\label{fig:fedllm_arch}
\end{figure}


\subsection{M3: TripleE-TGNN for Multi-Level Security Monitoring}

The service graph $\calG_s(t)$ has nodes representing microservices and edges representing API calls. The trace graph $\calG_r(t)$ has nodes representing distributed traces. The host graph $\calG_n(t)$ has nodes representing infrastructure hosts. EWC consolidation strength $\lambda_{\mathrm{ewc}}=0.5$ balances plasticity and stability during topology evolution. Multi-task loss weights are $\lambda_s=0.4$, $\lambda_r=0.3$, $\lambda_n=0.3$ with consistency weight $\lambda_{\mathrm{con}}=0.1$.

\begin{figure}[H]
\centering
\includegraphics[width=0.75\textwidth]{figures/fig_triplee_tgnn_architecture.pdf}
\caption{TripleE-TGNN architecture for multi-level security monitoring in cloud-native environments.}
\label{fig:triplee_arch}
\end{figure}


\subsection{M4: MambaShield for Drift-Adaptive Detection}

MambaShield~\cite{Anaedevha2026mamba} is instantiated for drift-adaptive detection. Sequential flow records are processed by the selective state-space model whose input-dependent transitions adapt to the current traffic regime. The state dimension is $d_s=64$, discretisation step $\Delta=0.01$, and FFT convolution kernel length $L=1024$. The teacher ensemble comprises $M=5$ models trained on disjoint data subsets. The curriculum schedule starts at $\rho_0=0.05$ poisoning fraction and increases to $\rho_{\max}=0.3$. PAC-Bayesian prior variance $\sigma_P^2=1.0$.

\begin{figure}[H]
\centering
\includegraphics[width=0.75\textwidth]{figures/fig_mambashield_architecture.pdf}
\caption{MambaShield architecture for drift-adaptive intrusion detection with progressive robustness distillation.}
\label{fig:mamba_arch}
\end{figure}


\subsection{M5/M6: Uncertainty-Calibrated Threat Assessment}

The stochastic transformer uses $M=50$ Monte Carlo forward passes at inference to produce epistemic and aleatoric uncertainty estimates. Predictions with epistemic uncertainty exceeding threshold $\tau=0.3$ are routed to human analysts, reducing alert fatigue by 43\%. The UC-HGP module provides hierarchical Gaussian process uncertainty at time scales matching the CT-TGNN branches ($\tau_1$, $\tau_2$, $\tau_3$). Variational learning rate $\alpha_\phi=5\times10^{-4}$; perturbation learning rate $\alpha_\psi=10^{-3}$.

\begin{figure}[H]
\centering
\includegraphics[width=0.96\textwidth]{figures/fig_stochastic_transformer_architecture.pdf}
\caption{Stochastic transformer with uncertainty-driven triage for security operations.}
\label{fig:stochastic_arch}
\end{figure}


\subsection{M7: Game-Theoretic Framework for Adaptive Adversaries}

The defender's action space $\calC$ comprises detection configurations: monitoring coverage, alert thresholds, and response policies. The adversary's action space $\calA$ encompasses target selection, perturbation type, and evasion strategy. The multiplicative weights learning rate $\eta=\sqrt{\log|\calC|/T}$. Uncertainty penalty $\kappa=0.5$; expert-referral threshold $\tau=0.3$.

\begin{figure}[H]
\centering
\includegraphics[width=0.75\textwidth]{figures/fig_game_theoretical_architecture.pdf}
\caption{Game-theoretic framework for adaptive graph adversaries in network security.}
\label{fig:game_arch}
\end{figure}


\subsection{Forward-Compatible Detection for Post-Quantum Protocols}

Protocol-specific encoders~\cite{Anaedevha2026edm,NIST2024,Regev2009} process timing features for lattice-based protocols (Kyber/Dilithium~\cite{Avanzi2022}), structural features for code-based protocols (McEliece/Niederreiter), and hash-based features (SPHINCS+). Cross-attention fusion combines the three representations. Adversarial training budget $\epsilon=0.03$ with PGD-40 perturbations.

\begin{figure}[H]
\centering
\includegraphics[width=0.96\textwidth]{figures/fig_pq_idps_architecture.pdf}
\caption{Forward-compatible multi-protocol detection architecture for post-quantum traffic.}
\label{fig:pqidps_arch}
\end{figure}


%%=====================================================================
\section{Hyperparameter Configuration and Optimisation}
\label{sec:hyperparams}
%%=====================================================================

Table~\ref{tab:hyperparams} consolidates the principal hyperparameters for all methods as configured for the network-security domain. All values were selected through grid search on the CIC-IoT-2023 validation set and confirmed stable across the remaining five datasets.

\begin{table}[htbp]
\centering
\caption{Consolidated hyperparameter configuration for all methods in the cybersecurity domain.}
\label{tab:hyperparams}
\small
\begin{tabular}{p{3.5cm} p{2.5cm} p{2.5cm} p{5.0cm}}
\toprule
\textbf{Method} & \textbf{Parameter} & \textbf{Value} & \textbf{Rationale} \\
\midrule
\multirow{4}{3.5cm}{M1: CT-TGNN}
& Hidden dim & 256 & Capacity for 83-feature input \\
& ODE tolerance & $10^{-5}$ & Balance accuracy/speed \\
& Time constants & 1/60/3600\,s & Packet/session/campaign scales \\
& Lipschitz bound & Spectral norm & Certified robustness \\
\addlinespace
\multirow{4}{3.5cm}{M2: FedLLM-API}
& Clipping $C$ & 1.0 & DP sensitivity bound \\
& DP $(\epsilon,\delta)$ & $(0.85, 10^{-5})$ & Strong privacy guarantee \\
& OT reg $\lambda$ & 0.1 & Sinkhorn convergence \\
& LLM weight & 0.3 & Ensemble balance \\
\addlinespace
\multirow{3}{3.5cm}{M3: TripleE-TGNN}
& EWC $\lambda$ & 0.5 & Forgetting prevention \\
& Level weights & 0.4/0.3/0.3 & Service emphasis \\
& Consistency $\lambda_{\mathrm{con}}$ & 0.1 & Cross-level agreement \\
\addlinespace
\multirow{3}{3.5cm}{M4: MambaShield}
& State dim & 64 & SSM capacity \\
& Kernel length & 1024 & Sequence coverage \\
& Poison schedule & 0.05--0.3 & Progressive hardening \\
\addlinespace
\multirow{3}{3.5cm}{M5/M6: Stoch.Trans.}
& MC samples & 50 & Uncertainty quality \\
& Triage threshold & 0.3 & FP/FN balance \\
& Calibration & Temperature & ECE minimisation \\
\addlinespace
\multirow{2}{3.5cm}{M7: Certification}
& MWU $\eta$ & $\sqrt{\log|\calC|/T}$ & Optimal regret \\
& Unc. penalty $\kappa$ & 0.5 & Risk sensitivity \\
\bottomrule
\end{tabular}
\end{table}

All methods are implemented in PyTorch~2.1~\cite{Paszke2019} with CUDA~12.1. Training uses AdamW with initial learning rate $2\times10^{-4}$, weight decay $10^{-5}$, and cosine annealing. Mixed-precision training reduces memory by 40\%. Gradient clipping at norm 1.0 prevents instability. Early stopping with patience of 10 epochs monitors validation loss.


%%=====================================================================
\section{Integrated System Architecture}
\label{sec:integrated_arch}
%%=====================================================================

The seven methods (M1--M7) are integrated into a unified detection pipeline whose data flow proceeds as follows. Raw network traffic enters through a data ingestion module that performs protocol parsing, flow extraction, and 83-feature computation. The CT-TGNN module constructs the temporal graph and computes continuous-time node embeddings. The TripleE-TGNN module maintains multi-level embeddings for hierarchical environments. The MambaShield module processes sequential event records for drift-adaptive scoring. The Stochastic Transformer and UC-HGP modules provide calibrated uncertainty estimates. All outputs feed into a decision engine that applies the game-theoretic strategy from the M7 Certification framework, producing alerts classified by severity and confidence. The FedLLM-API module enables periodic model synchronisation across organisational boundaries.

\begin{figure}[H]
\centering
\includegraphics[width=0.96\textwidth]{figures/Main_arch_v1b.pdf}
\caption{Integrated detection pipeline showing the data flow from raw network traffic through all seven methods to severity-ranked, confidence-scored alerts. The FedLLM-API module synchronises models across boundaries under differential privacy.}
\label{fig:u_idps_arch}
\end{figure}

The system is designed for deployment as a plug-in for existing IDPS frameworks such as Snort~3~\cite{Roesch1999} and Suricata, consuming their alert feeds as additional input and producing enriched, confidence-scored outputs. The full software implementation is described in Chapter~\ref{ch:platform}, publicly archived under DOI 10.5281/zenodo.19129512~\cite{Anaedevha2026robustidps}.


%%=====================================================================
\section{Chapter Summary}
\label{sec:ch3_summary}
%%=====================================================================

This chapter has demonstrated how the domain-independent mathematical methods of Chapter~\ref{ch:methods} are mapped onto the entities and constraints of network intrusion detection. The key contributions of this chapter are: (1)~a formal data-to-graph transformation pipeline converting raw PCAP/CSV network traffic into temporal graph structures with node features (12 dimensions) and edge features (83 dimensions); (2)~a unified classification taxonomy of 34 classes across 7 attack categories, extensible through zero-shot LLM classification; (3)~a comprehensive 83-feature engineering pipeline spanning flow statistics, inter-arrival distributions, flag indicators, payload statistics, and derived ratios; (4)~domain-specific instantiation of all seven methods with cybersecurity-tailored hyperparameters, time constants, and architectural configurations; and (5)~an integrated detection pipeline demonstrating the end-to-end flow from raw traffic to confidence-scored alerts. The next chapter presents the complete experimental programme validating these domain-specific instantiations on six benchmark datasets totalling 84.2 million labelled samples.

